\documentclass[11pt]{article}
\usepackage[margin=1in]{geometry}
\usepackage{authblk}
\usepackage{graphicx}
\usepackage{booktabs}
\usepackage{amsmath, amssymb}
\usepackage{caption}
\captionsetup{font=small,labelfont=bf}
\usepackage{array}
\usepackage[hidelinks]{hyperref}

\title{Trademark filings and declining US business dynamism:\\
coverage of business formation, three failed tests of the\\
crowding hypothesis, and surname-based ethnic composition}
\author{Eric Silver}
\affil{}
\date{May 2026}

\begin{document}
\maketitle

\begin{abstract}
We measure what fraction of US new-business formation is captured by
USPTO trademark filings, by joining the annual count of \emph{first-ever
trademark filers} (each owner's earliest filing across all 45 NICE
classes) to the Census Bureau's Business Formation Statistics (BFS),
2004--2024. USPTO debut filers cover \textbf{$5$--$8\%$ of all SS-4
business applications (BA)} and \textbf{$7$--$22\%$ of High-Propensity
Business Applications (HBA, the employer-likely subset)}, with the
HBA coverage ratio rising secularly from $\sim 7\%$ in 2005 to
$\sim 18\%$ in 2024 (with a pandemic-era spike to $22\%$ in 2020). This
brackets the $10$--$15\%$ employer-business coverage estimate reported
by Dinlersoz, Goldschlag, Myers, \& Zolas (2018) using restricted
Census Longitudinal Business Database (LBD) microdata. We extend
their result by documenting (i) the upward trend in coverage over
two decades, indicating that new businesses are progressively more
trademark-aware, and (ii) by applying the 2010 Census surname /
race-ethnicity probability file (Word, Coleman, Nunziata, \& Kominski,
2008) to the $1.83$M individual-name debut filers in our corpus. The
surname-based composition of individual filers shifts dramatically
toward Asian/Pacific-Islander-classified names ($5\%$ in the 1990s,
$33\%$ in 2020--2024); we flag the well-documented $2018$--$2021$
foreign-applicant filing surge from China as a confound and discuss
what an unconfounded estimate would require (applicant country from
the USPTO XML, not currently in our processed panel).

\textbf{Three failed tests of the trademark-crowding hypothesis.}
Motivated by the declining-business-dynamism literature
(Decker et al.\ JEP, 2014; Akcigit \& Ates AEJ:Macro, 2021) and the
``kill zones'' framework
(Kamepalli, Rajan, \& Zingales AEJ:Micro, 2022;
Cunningham, Ederer, \& Ma JPE, 2021), we test three predictions of
the hypothesis that incumbent firms are crowding out new-entrant
trademark activity. All three fail. First, the pooled
first-time-filer share is \emph{rising}, not falling: $46$--$49\%$
trough in $1996$--$2003$, $56$--$58\%$ at present. Second, the top-$10$
owner share of filings is \emph{falling} in every NICE class --- the
``FAANG concentration'' story is not visible in trademark filings.
Third, an IPO event study on $877$ class-events from $2{,}536$ public
firms finds no chilling effect on within-class first-time-filer
counts at $t\pm 3$ years around the IPO (post coefficient
$+0.007$ on $\log(1 + \text{debuts})$, $t = +0.67$, $p = 0.50$).
The conventional crowding story that would tie declining business
dynamism to trademark concentration is not supported. The selective
within-class signal that does survive --- falling debut share in
tech-info classes (Sci~\&~Tech $042$, Telecommunications $038$) ---
is driven by mid-tier multi-filer proliferation (not top-$10$
concentration) and even that fails as a causal mechanism in the
within-class panel regression with class and year fixed effects.
\end{abstract}

\section{Why this measurement question matters}

\paragraph{Declining business dynamism.} The US start-up rate has
fallen from $\sim 12\%$ of firms per year in the late 1980s to
$\sim 8\%$ in the 2010s, with the decline visible across nearly every
sector (Decker, Haltiwanger, Jarmin, \& Miranda, 2014, \emph{JEP};
Akcigit \& Ates, 2021, \emph{AEJ: Macro}). Several explanations are on
the table --- demographic slowing of population growth (Hopenhayn,
Neira, \& Singhania, 2022, \emph{Econometrica}); rising fixed costs of
entry; increasing industry concentration --- but the data
infrastructure for testing them is incomplete: the gold-standard
Census Business Dynamics Statistics (BDS) is lagged by $\sim 3$ years
and anonymises owner names; the BFS series goes back only to $2004$
and tracks EIN applications (which include estates and household
employers, not just business starts).

\paragraph{Trademarks as a complementary signal.} The USPTO trademark
corpus offers something neither BDS nor BFS provides: a near-real-time,
name-disclosed register of commercial activity covering nearly every
sector of the consumer-facing economy, with classifications (NICE),
goods/services prose, and direct linkage to identifiable applicants.
The question this paper asks is the prerequisite for using trademark
data on the dynamism question: \emph{how much of new-business activity
do trademark debuts actually capture?}

\section{Coverage of US business formation, 2004--2024}

\subsection*{Method}
For every owner name appearing in the clean H=2y panel of
$10{,}060{,}073$ filings (1990--2021 cohorts; methodology in
\texttt{paper/main.tex}), we identify the owner's earliest filing
across all NICE classes and call that filing the owner's
\emph{debut}. We count unique-owner debuts by year. We compare to the
Census BFS time series at
\texttt{census.gov/econ/bfs/csv/bfs\_monthly.csv}, summing 12 months
of \emph{unadjusted} counts to an annual total. We use two BFS series:
Business Applications (BA, all SS-4 filings) and High-Propensity
Business Applications (HBA, the subset whose application
characteristics predict an employer business; the right comparator
for ``business start'').

\begin{figure}[h]
\centering
\includegraphics[width=\linewidth]{results/dynamism_coverage.png}
\caption{\textbf{USPTO trademark debut filers vs.\ Census BFS,
2004--2024.} Top: annual counts in millions of filings.  Bottom:
coverage ratio. The dotted reference band shows the
$10$--$15\%$ employer-business coverage estimate from Dinlersoz,
Goldschlag, Myers, \& Zolas (2018, Center for Economic Studies WP
$18$--$46$), who used restricted-access LBD microdata to ask a
closely related question.}
\label{fig:coverage}
\end{figure}

\subsection*{Headline numbers}

\begin{table}[h]
\centering
\small
\begin{tabular}{lrrrrr}
\toprule
Period & USPTO debuts & BFS BA & BFS HBA & USPTO/BA & USPTO/HBA \\
\midrule
2004--2008 &   $464{,}273$ & $10{,}293{,}896$ & $5{,}768{,}077$ & $4.5\%$ & $8.0\%$  \\
2014--2018 &   $907{,}851$ & $15{,}104{,}778$ & $6{,}199{,}297$ & $6.0\%$ & $14.6\%$ \\
2019--2024 & $1{,}753{,}741$ & $29{,}002{,}345$ & $9{,}963{,}498$ & $6.0\%$ & $17.6\%$ \\
\bottomrule
\end{tabular}
\caption{USPTO debut filers as a share of Census BFS business
applications, three five-year bands. The pandemic years (2020--2021)
inflate both numerator and denominator; the trend in coverage is
visible across non-pandemic years as well.}
\label{tab:coverage}
\end{table}

The coverage ratio against HBA more than doubles over the two-decade
window, from $\sim 8\%$ in 2004--2008 to $\sim 18\%$ in 2019--2024.
Three interpretive readings:

(1) \textbf{Tilt of new-business formation toward IP-intensive
sectors.} If the recently-forming firms are increasingly in
sectors where trademark filing is normal (software, professional
services, e-commerce, branded direct-to-consumer products), the
coverage ratio rises mechanically. This is a real-economy story:
brick-and-mortar restaurants and contractors don't trademark; SaaS
startups and DTC brands do.

(2) \textbf{Falling cost of filing.} Online filing through TEAS Plus
and TEAS Standard (in production since the mid-2000s) reduced the
effort and expense of trademark filing. The same firm-type in 2024
files at a higher rate than it would have in 2005.

(3) \textbf{Selection on quality.} If the BFS denominator captures
many ``shell'' or non-operating SS-4 filings (estates, holding
companies, household employers), the rising ratio reflects a
constant operating-firm-trademark rate against a denominator
increasingly diluted by non-operating EIN registrations.

The three are not mutually exclusive. Distinguishing them requires
sectoral decomposition, which we cannot do at present because our
processed parquets do not retain NAICS-equivalent applicant industry
codes and the NICE-NAICS crosswalk is imperfect (Hall, Helmers,
Rogers, \& Sena, 2014, \emph{Research Policy}). Sectoral
decomposition is a clear next step.

\subsection*{Comparison to Dinlersoz et al.\ (2018)}

Dinlersoz, Goldschlag, Myers, \& Zolas (2018), ``An Anatomy of US
Firms Seeking Trademark Registration,'' Center for Economic Studies
working paper $18$--$46$, ran the analogous accounting inside Census
using restricted-access LBD microdata, asking what fraction of new
\emph{employer} businesses ever file a trademark. They reported a
range of $\sim 10$--$15\%$ depending on sector and definition. Our
public-data result of $7$--$22\%$ against HBA spans their range and
extends it through 2024 with a clear upward trend. The mid-decade
overlap is close: 2014--2018 in our data is $14.6\%$, very near
their $10$--$15\%$ point estimate for the same time period.

\section{Surname-based race/ethnicity composition}

\subsection*{Method}
The Census Bureau's 2010 surname file (Word et al., 2008) reports,
for each surname appearing $\geq 100$ times in the 2010 Census, the
percentage probability that a person with that surname identifies as
each of: non-Hispanic White, non-Hispanic Black, Asian/Pacific
Islander, American Indian/Alaska Native, Two or More Races, or
Hispanic (any race). The file covers $162{,}254$ surnames.

For each USPTO debut owner whose name parses as a personal name
(2--4 alphabetic tokens, no entity markers such as
\texttt{LLC}/\texttt{INC}/\texttt{CORP}, no digits/URLs/ampersands),
we extract the candidate surname (the last token after stripping
suffixes \texttt{JR}, \texttt{SR}, \texttt{II}, \texttt{MD}, etc.)
and look up its Census probabilities. Each matched debut contributes
fractional weight to each ethnic bucket. Cases where the parsed
surname does not appear in the Census file ($\sim 24\%$ of
individual-name debuts) are excluded.

\paragraph{Sample sizes.}
$5{,}674{,}909$ unique debut owners total. $1{,}827{,}163$ parse as
individual personal names ($32.2\%$ of debuts; the rest are entities).
$1{,}390{,}267$ of those individual names match a surname in the
Census file ($76.1\%$ match rate). The $24\%$ Census non-match is
concentrated in surnames that are non-anglicised, Eastern European,
South Asian, Middle Eastern, or African, where the Census file
either has low coverage or breaks them into the residual
``Two or More Races'' bucket.

\subsection*{Results}

\begin{figure}[h]
\centering
\includegraphics[width=\linewidth]{results/surname_ethnicity.png}
\caption{\textbf{Estimated race/ethnicity composition of USPTO
individual-name trademark debut filers, 1990--2024,} via Census 2010
surname probabilities (Word, Coleman, Nunziata, \& Kominski, 2008).
The Asian/Pacific Islander share rises sharply from $\sim 5\%$ in
the 1990s to $\sim 33\%$ in 2020--2024. Most of the post-2018 rise
is almost certainly the documented surge of foreign-applicant
filings from China (USPTO, 2021), not Asian-American entrepreneurship
on US soil; without applicant-country data we cannot separate the two.}
\label{fig:surname}
\end{figure}

\begin{table}[h]
\centering
\small
\begin{tabular}{lrrrrrr}
\toprule
Period & White & Black & API & Hispanic & 2+\,races & AIAN \\
\midrule
1990--1999 & $74.6\%$ & $11.0\%$ & $5.2\%$  & $6.3\%$  & $1.8\%$ & $0.6\%$ \\
2000--2009 & $70.3\%$ & $11.6\%$ & $6.9\%$  & $8.4\%$  & $1.8\%$ & $0.6\%$ \\
2010--2019 & $61.4\%$ & $11.7\%$ & $14.2\%$ & $9.9\%$  & $1.8\%$ & $0.6\%$ \\
2020--2024 & $43.9\%$ & $10.7\%$ & $33.4\%$ & $9.4\%$  & $1.8\%$ & $0.5\%$ \\
\bottomrule
\end{tabular}
\caption{Surname-imputed race/ethnicity shares among USPTO debut
filers whose owner name parses as a personal name and whose surname
matches the 2010 Census file. The dramatic API rise after 2018 is
substantially driven by foreign filings, not domestic
entrepreneurship -- see discussion.}
\label{tab:surname}
\end{table}

\paragraph{The post-2018 API rise is mostly foreign filings.} The
Asian/Pacific Islander share goes from $\sim 7\%$ in 2009 to $\sim 18\%$
by 2018 to $\sim 33\%$ in 2020--2024. The largest single component of
the post-2018 surge is well-documented: USPTO publicly reported a
massive influx of trademark filings from Chinese applicants in
$2017$--$2021$, driven by subsidies from local Chinese governments
for filing US trademarks (USPTO, 2021; Hicks, 2024, \emph{Stanford
Law Review}). USPTO responded by requiring US-licensed counsel for
foreign-domiciled applicants ($2019$) and tightening signature
verification ($2021$). Without applicant-country data --- which the
USPTO XML includes but which our current parser does not extract ---
we cannot separate ``rise of Asian-American entrepreneurship in the
US'' from ``rise of Chinese-applicant filings under the US system.''
The pre-$2018$ rise from $5\%$ to $\sim 14\%$ likely reflects a
genuine demographic shift; the post-$2018$ rise of an additional
$20$ pp does not.

\paragraph{Other shares are more stable.} Non-Hispanic Black share
is remarkably flat at $\sim 11\%$ across the entire $1990$--$2024$
window, consistent with the absence of a major foreign-filing
confound from majority-Black countries. Hispanic share rises modestly
from $\sim 6\%$ in the 1990s to $\sim 10\%$ in the 2010s and plateaus.
Non-Hispanic White share falls from $\sim 75\%$ to $\sim 44\%$ --- but
the bulk of that fall is the API rise's mirror image, so
disentangling foreign vs.\ domestic API is required to interpret it.

\paragraph{The selection problem.} Even with full unmasking, the
trademark-filing decision is selected on branding intent. Hispanic-owned
small businesses (taquerias, contractors, landscapers) trademark at
lower rates per business than tech founders do, so the Hispanic
share of trademark debuts \emph{underestimates} the Hispanic share
of business starts in absolute terms. The rising-over-time pattern is
more interpretable than the absolute level.

\section{Multi-filer mix and the first-time-filer share, 1985--2025}
\label{sec:firsttimer-share}

For each year $y$ and each NICE class $c$, we compute the number of
filings whose owner had no earlier trademark filing in any class ---
the \emph{debut} filings --- and the total filing count. The
debut share $s_{c,y} = n_\text{debut} / n_\text{total}$ measures the
share of class-year filing activity from first-time filers.

Figure~\ref{fig:firsttimer-pooled} plots the pooled (across all 44
NICE classes with $\geq 5{,}000$ filings) annual totals.
Total filings grew from $71{,}459$ in $1985$ to $832{,}505$ in $2025$
(a factor of $11.6\times$). Debut filings grew from $37{,}705$ to
$469{,}199$ (a factor of $12.4\times$, slightly faster). The pooled
debut share is \emph{not falling}: a trough at $46$--$49\%$ in
$1996$--$2003$, a rise to $52$--$56\%$ from $2015$ onward, and a
pandemic-era peak of $57$--$58\%$ in $2020$--$2025$. This trajectory
is the opposite of the falling new-entrant share that the conventional
declining-business-dynamism story would predict.

\begin{figure}[h]
\centering
\includegraphics[width=\linewidth]{results/dynamism_firsttimer_share_pooled.png}
\caption{Pooled USPTO trademark filing counts and first-time-filer
(debut) share, 1985--2025. Total filings (grey, left axis) and debut
filings (blue, left axis) both rise by $\sim 12\times$ over the
window; debut share (red dashed, right axis) is flat-to-rising, not
falling.}
\label{fig:firsttimer-pooled}
\end{figure}

\subsection*{Sectoral decomposition: where the trend hides}

The pooled rise in the debut share is composition-driven. We
decompose $\Delta s_\text{pooled} = \sum_c w_c\,\Delta s_c + \sum_c
s_c^{\text{lag}}\,\Delta w_c$ where $w_c$ is class $c$'s share of
total filings. The within-class effect ($w_c\,\Delta s_c$) and the
composition effect ($s_c\,\Delta w_c$) are reported separately in
Table~\ref{tab:firsttimer-decomp} for the $2001$--$2025$ rise.

\begin{table}[h]
\centering
\small
\begin{tabular}{lrr}
\toprule
Industry & within ($\Delta$ share) & composition ($\Delta$ weight) \\
\midrule
\textbf{Largest positive contributors to debut-share rise, 2001--2025} & & \\
  Clothing, Footwear \& Headwear & +0.44 pp & +1.88 pp \\
  Medical \& Personal Care & +0.20 pp & +1.44 pp \\
  Education \& Entertainment & +0.37 pp & +1.18 pp \\
  Household Utensils & +0.36 pp & +1.17 pp \\
  Cosmetics \& Cleaning & +0.59 pp & +0.71 pp \\
\midrule
\textbf{Largest negative contributors} & & \\
  Insurance \& Financial & +0.04 pp & -1.24 pp \\
  Telecommunications & -0.24 pp & -1.56 pp \\
  Software \& Electronics & +0.27 pp & -2.21 pp \\
  Advertising \& Retail & -0.09 pp & -2.03 pp \\
  Scientific \& Tech Services & -0.53 pp & -4.78 pp \\
\bottomrule
\end{tabular}

\caption{Top-5 and bottom-5 NICE-class contributions to the
$2001$--$2025$ pooled debut-share rise, decomposed into within-class
(change in class-specific debut share at fixed weight) and
composition (change in class's share of total filings at fixed
within-class debut share). The pooled $+10$ percentage-point rise is
roughly $35\%$ within-class and $65\%$ composition: classes with
high within-class debut shares (consumer goods, services) grew as a
share of total filings, while classes with lower within-class debut
shares (Sci~\&~Tech $042$, Software $009$, Advertising $035$) shrank.}
\label{tab:firsttimer-decomp}
\end{table}

The five largest positive contributors are
Clothing/Footwear (NICE $025$, $+2.3$ pp),
Medical \& Personal Care ($044$, $+1.6$ pp),
Education \& Entertainment ($041$, $+1.6$ pp),
Household Utensils ($021$, $+1.5$ pp), and
Cosmetics \& Cleaning ($003$, $+1.3$ pp).
The five largest negative contributors are
Scientific~\& Tech Services ($042$, $-5.3$ pp),
Advertising~\& Retail ($035$, $-2.1$ pp),
Software~\& Electronics ($009$, $-1.9$ pp),
Telecommunications ($038$, $-1.8$ pp), and
Insurance~\& Financial ($036$, $-1.2$ pp). The economy-wide debut
share rose mostly because services and consumer-goods classes
(where most filings are first-time) grew as a share of total
trademark activity, while tech-information classes shrank.

\subsection*{Within-class debut share: where it does fall}

Within tech-information classes specifically, the debut share is
falling. Three illustrative time-series slopes ($1985$--$2025$
linear fit to within-class debut share):
\begin{itemize}\setlength\itemsep{0pt}
\item Scientific \& Tech Services (NICE $042$): $-0.42$ pp/year
  (debut share $67.5\% \to 54.5\%$).
\item Telecommunications (NICE $038$): $-0.08$ pp/year
  (debut share $59.0\% \to 46.1\%$).
\item Software \& Electronics (NICE $009$): $+0.06$ pp/year, flat
  near $53\%$.
\end{itemize}
By contrast, consumer-services classes have rising within-class debut
shares: Medical~\& Personal Care ($044$) $51\% \to 64\%$,
Legal~\& Personal Services ($045$) $50\% \to 63\%$,
Restaurants~\& Hotels ($043$) $49\% \to 60\%$. The next two sections
ask whether the within-class debut-share decline in tech-information
classes is mechanically explained by incumbent owner concentration.

\section{Top-$N$ owner concentration: the FAANG-crowding hypothesis fails}
\label{sec:topN-concentration}

For each (class, year) we computed the top-$10$ owner share of
filings and the Herfindahl index on owner shares
(Figure~\ref{fig:topN-byclass}). \textbf{Every NICE class shows
top-10 owner share decreasing over the panel.} The most dramatic
declines are in Telecommunications ($038$, HHI $48.4 \to 4.4$),
Sci~\& Tech ($042$, HHI $3.3 \to 0.5$),
Software ($009$, HHI $3.2 \to 0.4$), and
Beverages ($032$, top-10 share $15.1\% \to 2.4\%$). Even in classes
where the within-class debut share is falling, the top-$10$
incumbents have a \emph{lower} share of filings than they did in
the 1990s.

\begin{figure}[h]
\centering
\includegraphics[width=\linewidth]{results/dynamism_topN_share_byclass.png}
\caption{Per-NICE-class top-10 owner share of filings, 1985--2025.
Red lines: class-specific top-10 share; grey: pooled. No class
shows a rising top-10 share. The very-high top-10 shares in the
early years (e.g., $90$--$94\%$ in services classes $043$, $044$,
$045$) reflect tiny early-year sample sizes; the trend over the
1995--2025 window is more reliable. The FAANG-crowding hypothesis
that ``Google/Amazon/Microsoft/Apple are squeezing out new-entrant
trademark activity'' is not supported in any class.}
\label{fig:topN-byclass}
\end{figure}

The natural reconciliation with the tech-class debut-share decline
is that the rise in non-debut filings within tech is coming from
\emph{mid-tier} multi-filers (owners with $2$--$9$ filings per year),
not top-$10$ incumbents. We test this directly.

\section{Mid-tier proliferation: the reframed hypothesis also fails}
\label{sec:midtier}

For each (class, year) we partition filings by owner intensity:
single-filer ($k = 1$ filing in that class-year), mid-tier
($2 \le k \le 9$), heavy ($k \ge 10$), and top-$10$. We then
regress the year-over-year change in debut share on the lagged
year-over-year changes in mid-tier, heavy, and top-$10$ shares,
with class and year fixed effects and standard errors clustered by
class:
$$\Delta s_{c,y}^{\text{debut}} \;=\; \beta_1 \Delta s_{c,y-1}^{\text{midtier}}
   \;+\; \beta_2 \Delta s_{c,y-1}^{\text{heavy}}
   \;+\; \beta_3 \Delta s_{c,y-1}^{\text{top10}}
   \;+\; \alpha_c + \alpha_y + \epsilon_{c,y}.$$

The crowding hypothesis predicts $\beta_1 < 0$ (mid-tier proliferation
absorbs filings that would otherwise be first-time). What we find
(Table~\ref{tab:midtier-regression}): $\beta_1 = +0.197$
($t = +1.95$, $p = 0.051$). The coefficient is the wrong sign.

\begin{table}[h]
\centering
\small
\begin{tabular}{lrrr}
\toprule
Predictor (lagged $\Delta$) & coef & cluster SE & $t$ \\
\midrule
  $\Delta$ mid-tier share (2--9 filings/year)   & +0.1965 & 0.1007 & +1.95 \\
  $\Delta$ heavy share (10+ filings/year)       & +0.1029 & 0.0807 & +1.27 \\
  $\Delta$ top-10 share                         & +0.0814 & 0.0399 & +2.04 \\
\midrule
  Class fixed effects & \multicolumn{3}{c}{Yes} \\
  Year fixed effects  & \multicolumn{3}{c}{Yes} \\
  $N$ (class $\times$ year cells) & \multicolumn{3}{c}{1,755} \\
  $R^2$  & \multicolumn{3}{c}{0.266} \\
\bottomrule
\end{tabular}
\caption{Within-class panel regression of $\Delta$debut share on
lagged $\Delta$multi-filer shares. All three coefficients are
positive --- the within-class year-over-year predictors all
co-move with debut share, not against it. Cluster-by-class
standard errors. $N = 1{,}755$ (class $\times$ year) cells.}
\label{tab:midtier-regression}
\end{table}

The substantive interpretation is that within a class, year-over-year
filing-volume shocks (e.g., the 2020 pandemic surge) raise both
mid-tier filing and first-time filing together: established firms
file more line-extension marks in busy years, and new firms file
their debut marks in those same busy years. Filing activity at
different intensity tiers \emph{co-moves} with the wave, not against
it. The mid-tier-proliferation crowding mechanism is not what
explains the tech-class within-class debut-share decline.

\section{IPO chilling at the NICE-class level: a clean null}
\label{sec:ipo-chilling}

The third test draws on the kill-zone literature
(Kamepalli, Rajan \& Zingales, 2022; Cunningham, Ederer \& Ma, 2021),
which finds that a Big-Tech acquisition or major IPO in a narrow
vertical chills venture-capital funding in adjacent verticals. The
analogous trademark question: when an IPO firm goes public, does
first-time-filer trademark activity in the IPO firm's NICE class
decline relative to the pre-IPO trend?

\paragraph{Event panel.} We join the USPTO--SEC owner crosswalk to
the SEC firm-year financials and define each public firm's
\emph{IPO year} as the first fiscal year in which it appears in our
SEC data, restricted to firms with first SEC year $\ge 2012$ (when
Census FSDS coverage stabilises) and USPTO first-trademark year
within $5$ years of the IPO (filtering out pre-existing public firms
that merely appear in the SEC data later). This yields $2{,}536$
IPO-firm events. For each IPO firm we identify the NICE classes in
which it actively files trademarks ($\ge 5$ filings in the
$\pm 2$ year window around IPO), producing $877$ active-filer
class-events spanning $230$ unique (class, IPO-year) cells.

\paragraph{Event-study specification.} For each class-event, we
stack non-IPO-firm debut counts at relative years $t \in [-3, +3]$
around the IPO. Two-way fixed-effects regression:
$$\log\bigl(1 + n_\text{debut}^{c,y}\bigr) = \beta\,\mathbf{1}\{y \ge y_e\}
   + \alpha_c + \alpha_y + \epsilon_{c,y,e}.$$
Cluster-by-class standard errors. We also fit a separately-sloped
specification, replacing the post indicator with two slopes
(pre-event and post-event) on the relative-year axis.

\paragraph{Result.} No detectable chilling effect
(Table~\ref{tab:ipo-chilling}). The post coefficient is $+0.007$
($t = +0.67$, $p = 0.50$); the implied multiplicative effect on
debut counts is $e^{+0.007} = 1.007$, indistinguishable from one.
The slope decomposition is also weak: pre-event slope
$+0.0042$/year, post-event slope $-0.0020$/year, difference
$-0.0063$/year. Figure~\ref{fig:ipo-chilling} shows the mean
trajectory of debut counts across event years $-3$ to $+3$.

\begin{table}[h]
\centering
\small
\begin{tabular}{lrrr}
\toprule
Estimator & coef & cluster SE & $t$ \\
\midrule
  $\mathbf{1}[t \geq \text{IPO}]$ (post vs pre) & +0.0069 & 0.0103 & +0.67 \\
  pre-event slope (rel\_neg) & +0.0042 & 0.0038 & +1.11 \\
  post-event slope (rel\_pos) & -0.0020 & 0.0025 & -0.82 \\
\midrule
  Class FE + year FE & \multicolumn{3}{c}{Yes} \\
  $N$ event-cells & \multicolumn{3}{c}{6,139} \\
  IPO-firm events & \multicolumn{3}{c}{2,536} \\
\bottomrule
\end{tabular}
\caption{Event-study DiD for the IPO-chilling test. Post coefficient
indistinguishable from zero at conventional significance, both as a
post-vs-pre indicator and as a slope difference. $877$ class-events
from $2{,}536$ IPO-firm events, $2012$--$2022$ IPO years.}
\label{tab:ipo-chilling}
\end{table}

\begin{figure}[h]
\centering
\includegraphics[width=\linewidth]{results/dynamism_ipo_chilling.png}
\caption{Mean log debut count by event-relative year for $877$
IPO class-events. The mean continues to rise across the event
window; the IPO at $t = 0$ does not produce a visible break in the
trajectory. Class-level test only --- the finer token-vertical
specification (cosine-similarity neighbourhoods of the IPO firm's
goods/services vocabulary) is left for future work.}
\label{fig:ipo-chilling}
\end{figure}

\paragraph{Scope of the null.} The class-level test is a coarse
operationalisation of the kill-zone hypothesis. The original VC kill
zone effect operates at the level of \emph{specific verticals}
defined by product-token overlap, not by SIC/NICE class. Our
$44$-NICE-class partition is coarser than the Kamepalli et al.\ unit
of analysis. We do not claim that no kill zones exist in
trademark-filing space; we report that the coarse class-level test
finds no effect, and flag the finer-grained token-vertical test
(cosine-similarity neighbourhoods of the IPO firm's pre-IPO
goods/services prose) as a follow-up that should be run before
declaring the hypothesis dead.

\section{What is left of the crowding story}

After three failed tests, the conventional incumbent-crowding
narrative does not survive at the trademark-filing level. What
\emph{is} visible in our data:
\begin{itemize}\setlength\itemsep{2pt}
\item The within-class debut-share decline in tech-information
  classes is a real $13$ pp drop (Sci~\&~Tech $042$:
  $67.5\% \to 54.5\%$) and a $13$ pp drop in Telecommunications
  ($038$). The drop is not explained by top-$10$ concentration
  (which fell) or mid-tier proliferation in a within-class panel
  sense.
\item The composition shift away from tech-information classes and
  toward consumer-services classes is itself the dominant driver of
  the rising pooled debut share. This is a sectoral story, not a
  concentration story.
\item Trademark activity is broadly \emph{decentralising}, not
  concentrating, across the entire $1985$--$2025$ panel. This
  contradicts the qualitative version of the declining-business-
  dynamism narrative even as it tracks the BFS quantitative trend
  in business applications.
\end{itemize}
The reconciliation we read in the data is that the BDS-measured
declining startup rate
(Decker et al., 2014; Akcigit \& Ates, 2021) is about
\emph{employer} firms (establishments with payrolls), while
trademark filings include large numbers of brand-aware solo
operators, DTC consumer-goods startups, and online businesses that
either do not employ or report through different establishment
structures. The two populations are diverging: employer business
formation is becoming less dynamic, brand-aware new-firm formation
is becoming \emph{more} dynamic, and the trademark corpus mostly
captures the latter.

\section{Next steps required for publication}

The current results are publishable as descriptive statistics with
the caveats above, but a complete treatment requires:

(1) \textbf{Re-parse the USPTO XML to extract applicant state and
country.} Our existing parser drops the \texttt{<address>} elements
under \texttt{<case-file-owner>}. With state, we can repeat the
coverage analysis at the state level (cross-cutting with BFS state
panels), and disentangle delivery-state vs.\ Wyoming/Delaware/Nevada
LLC-haven distortions. With country, we can separate foreign
applicants and produce an unconfounded API time series.

(2) \textbf{NAICS-NICE crosswalk for sectoral decomposition.}
Three candidate readings of the rising HBA coverage ratio (sectoral
tilt, falling filing cost, denominator dilution) are distinguishable
with a sectoral breakdown. The Hall et al.\ (2014, \emph{Research
Policy}) crosswalk is the standard reference.

(3) \textbf{Comparison panel with SBA PPP disclosures (2020--2021).}
PPP loan disclosures over \$150K named the owner explicitly,
covering $\sim 11$M small businesses including those that never
trademark. This is a one-time but very large unmasking opportunity
for the immigrant-business question specifically.

(4) \textbf{Validation against state Secretary of State new-entity
filings.} State-level cross-check of the BFS denominator using
California, Texas, Florida, and New York monthly entity-formation
counts.

(5) \textbf{Token-vertical IPO chilling test.} Replace the
NICE-class operationalisation in Section~\ref{sec:ipo-chilling}
with cosine-similarity neighbourhoods of the IPO firm's pre-IPO
goods/services prose (the same infrastructure as Section~5 of the
companion $\Delta KL$ paper). Compare core vertical (cosine
$\ge 0.7$), adjacent vertical (cosine $0.5$--$0.7$), and matched
non-event control verticals on debut-count trajectory. The
class-level null in this paper does not rule out a token-vertical
kill-zone effect.

(6) \textbf{Acquisition events as the canonical kill-zone test.}
The Kamepalli et al.\ result is on acquisitions, not IPOs.
Augmenting our event panel with Crunchbase or SEC 8-K-parsed
acquisitions ($2010$--$2024$) would let us replicate the kill-zone
test in its original form on trademark debuts.

(7) \textbf{Reconciling BDS vs.\ BFS vs.\ trademark debuts.}
The three series tell different stories about the same underlying
question. BDS (employer firm entry rate) is falling. BFS (all
business applications, BA) is rising sharply post-$2020$. Trademark
debuts are rising. A clean joint reconciliation requires
employer/non-employer disaggregation of trademark debuts using the
Census restricted LBD --- which is precisely what Dinlersoz et al.\
(2018) had access to. The version of their analysis with a
$2024$ cutoff is the natural next data ask.

\section*{Data and code}
All code is in the \texttt{business-dynamism} branch of the
project repository:
\texttt{scripts/business\_dynamism\_coverage.py} (coverage),
\texttt{scripts/surname\_ethnicity.py} (surname race/ethnicity),
\texttt{scripts/dynamism\_firsttimer\_share.py} (debut-share panel),
\texttt{scripts/dynamism\_class\_contributions.py} (decomposition),
\texttt{scripts/dynamism\_topN\_share.py} (top-$N$ concentration),
\texttt{scripts/dynamism\_midtier\_regression.py} (panel regression),
and
\texttt{scripts/dynamism\_ipo\_chilling.py} (event study).
BFS data is at \texttt{census.gov/econ/bfs/csv/bfs\_monthly.csv}.
The Census surname file is at
\texttt{www2.census.gov/topics/genealogy/2010surnames/names.zip}.

\section*{References}
\begin{itemize}\setlength\itemsep{2pt}
\item Akcigit, U., \& Ates, S.~T. (2021). Ten facts on declining
  business dynamism and lessons from endogenous growth theory.
  \emph{American Economic Journal: Macroeconomics}, $13$(1), $257$--$298$.
\item Cunningham, C., Ederer, F., \& Ma, S. (2021). Killer
  acquisitions. \emph{Journal of Political Economy}, $129$(3),
  $649$--$702$.
\item Kamepalli, S.~K., Rajan, R., \& Zingales, L. (2022). Kill
  zone. \emph{American Economic Journal: Microeconomics},
  forthcoming.
\item Decker, R.~A., Haltiwanger, J., Jarmin, R.~S., \& Miranda, J.
  (2014). The role of entrepreneurship in US job creation and
  economic dynamism. \emph{Journal of Economic Perspectives},
  $28$(3), $3$--$24$.
\item Dinlersoz, E., Goldschlag, N., Myers, A., \& Zolas, N. (2018).
  An anatomy of US firms seeking trademark registration. Center for
  Economic Studies Working Paper $18$--$46$. US Census Bureau.
\item Hall, B.~H., Helmers, C., Rogers, M., \& Sena, V. (2014). The
  choice between formal and informal intellectual property: A review.
  \emph{Journal of Economic Literature}, $52$(2), $375$--$423$.
\item Hopenhayn, H., Neira, J., \& Singhania, R. (2022). From
  population growth to firm demographics. \emph{Econometrica},
  $90$(4), $1879$--$1914$.
\item USPTO. (2021). Trademark Public Advisory Committee Annual
  Report. United States Patent and Trademark Office.
\item Word, D.~L., Coleman, C.~D., Nunziata, R., \& Kominski, R.
  (2008). Demographic aspects of surnames from Census 2000. US
  Census Bureau working paper.
\end{itemize}

\end{document}
